\section{Introduction}
\label{sec:introduction}

Legged robots can cross terrain that is broken, narrow, or hard for
wheels and tracks. They only need a sequence of usable footholds. The
hard part is choosing those footholds. The robot must also choose
which leg moves, when the leg leaves the ground, and how long the
swing lasts. This full problem is called \textit{contact planning}.
In this paper, a planner is \textit{greedy} if it commits only the
next immediate action. It has no explicit lookahead. We use
\textit{acyclic} for contact sequences that are not a periodic gait
cycle and are not transitions through a fixed gait library. The
literature also uses \textit{non-gaited} and \textit{free-gait} for
similar ideas.

\subsection{Related Work}

Classical locomotion planners often place foothold choice, contact
timing, and whole-body motion inside Model Predictive Control (MPC)
\cite{Kim2019highly-dynamic, Wensing2022Nov}. This works well when
the contact pattern is known. Once the planner must choose contacts,
the problem becomes a Mixed-Integer Quadratic Program
\cite{Geisert2019contact-planning}. Some methods reduce the search by
limiting footholds to convex terrain regions
\cite{deits_footstep_2014}. Others optimize contact schedules and
timing together \cite{winkler_gait_2018,
aceituno_simultaneous_2019}. These methods shrink the problem, but
they are still costly for fast reactive correction.
Other planners remove the periodic schedule without using a
learned policy. They choose contact timing online from per-leg utility
or feasibility rules
\cite{fankhauser_free_gait_2016, boussema_online_2019,
bjelonic_whole-body_2021}. Search-based non-gaited planners
\cite{amatucci_monte_2022} also make truly acyclic contact sequences.
Recent work speeds up these searches by learning the expensive part,
such as cost-to-go estimation \cite{taouil_non-gaited_2024} or
transition feasibility \cite{akizhanov_learning_2024}.
Contact-implicit methods avoid the discrete branch by relaxing
complementarity constraints. They can discover contacts from dynamics
in real time on hardware \cite{kim_contact-implicit_2025}. They do
not, however, choose footholds from a perceived steppability map. That
is the decision learned in this work.

End-to-end policies put the same discrete choices inside one
neural controller. They can cross hard terrain reliably
\cite{zhang_learning_2023, zhang_novel_2024, duan_sim--real_2022,
siekmann_blind_2021, lee_learning_2020}. But the contact decision is
hidden inside the joint targets. It cannot be checked or constrained
before execution. Perceptive model-based stacks take a different path.
They project a commanded foothold onto a convex terrain region before
tracking \cite{grandia_perceptive_2022, leggedperceptive}. This
requires an explicit foothold to project.

Hybrid methods keep learning more focused. A common design learns
foothold placement. A lower-level controller or gait generator still
provides the contact schedule
\cite{gangapurwala_rloc_2022, villarreal_fast_2019, omar_fast_2022,
omar_safesteps_2023, asselmeier_steppability-informed_2024,
gaspard_footstepnet_2024, peng_deeploco_2017, tsounis_deepgait_2020,
shi_terrain-aware_2023, xie_glide_2023}. Another design
learns contact structure and leaves foothold targets to a separate
heuristic. Examples include selecting whole-robot contact
configurations \cite{da_learning_2020}, phase parameters
\cite{yang_fast_2021}, or members of an offline gait library
\cite{sun_online_2024}. These methods can adapt timing or contact
mode. Still, their possible contact sequences are limited by a
structure chosen before deployment. Patrizi et al.
\cite{patrizi_rl-augmented_2026} loosen this structure by learning
per-foot flight-phase injections into an MPC horizon. This gives
acyclic timing changes, but it does not learn foothold selection.
End-to-end policies can also show emergent, speed-dependent gait
structure when trained without an explicit schedule
\cite{fu_minimizing_2021}.

The closest comparison is ContactNet \cite{bratta_contactnet_2024}.
It uses a neural network to rank discretized foothold candidates. It
does not prescribe which leg should swing. Its concurrency level is
set by configuration. In walking experiments, one leg swings at a
time. In trot experiments, it ranks joint footholds for the two legs
that swing together. It re-ranks candidates at touchdown. The
configured gait sets the swing timing. GaitNet differs in
three ways. It re-ranks during swing, selects swing duration, and lets
leg overlap emerge from repeated single-leg decisions instead of a
configured contact pattern.

\subsection{Contributions}

GaitNet makes one atomic decision at a fixed control rate. It chooses
which leg to swing, where to place it, and how long the swing lasts.
This decision can happen while another leg is already in the air. The
planner does not represent phase offsets, gait primitives, leg pairs,
or multi-leg actions. At each tick, it scores terrain-valid candidates
for all legs. It compares them with a learned no-action option. It
then commits at most one footstep for the whole robot. Any swing
overlap is therefore produced by the value function and the repeated
decision rule. Bratta et al. list re-ranking during swing as future
work \cite{bratta_contactnet_2024}. This paper evaluates that
mechanism directly. This narrow decision is useful because the
foothold remains explicit. The terrain projection step can still act
as a safety layer. The same frozen planner is also used with two
different low-level controllers: a convex MPC in simulation and an
NMPC$+$WBC stack on hardware.

The main contributions are:
\begin{enumerate}
  \item A greedy, high-rate footstep planner that emits at most one
    action per control tick. It produces multi-leg swing without an
    explicit phase schedule or a separate footstep-evaluation stage.
  \item A per-candidate formulation in which the learned
    network does not depend on the number of candidates. A policy
    trained on sampled candidates can therefore be deployed, frozen,
    on an exhaustive terrain-mask enumeration.
  \item Hardware validation on a Unitree Go1 against the Legged
    Perceptive Stack, a perceptive Nonlinear Model Predictive Control
    (NMPC) baseline. We also test concurrency-restricted GaitNet
    ablations to isolate the value of emergent multi-leg swing.
\end{enumerate}

This paper extends a preliminary framework introduced in previous
work \cite{Sullivan2025}. That work introduced the single-footstep,
single-tick action space.
It reported simulation results with randomly sampled candidate
footholds. Here, deployment replaces random sampling with
exhaustive enumeration over the same terrain mask. We also provide
the first hardware evaluation of the planner. We include preliminary
simulation results for context in \autoref{sec:results}. They are
marked as preliminary throughout. They were not rerun with the
deployment-time enumeration procedure used here.
